%! Minimizing Loss by Gradient Descent — Unit 8, Lesson 8.3 — Student Packet Cover Sheet
\documentclass[10pt]{article}
\usepackage{linalg-article}
\usepackage{linalg-boxes}
\usepackage{ltablex}
\keepXColumns

\begin{document}

% Full-bleed burgundy banner
\begin{tikzpicture}[remember picture, overlay]
  \fill[burgundy] (current page.north west)
    rectangle ([yshift=-0.9in]current page.north east);
\end{tikzpicture}
\vspace{-0.8in}

\begin{center}
  {\color{white}\textbf{\LARGE Linear Algebra}}\\[4pt]
  {\color{blushmid}\large\textbf{Unit 8 \quad Learning from Data}}\\[6pt]
  {\color{white}\textbf{\large Lesson 8.3 \quad Minimizing Loss by Gradient Descent}}
\end{center}

\vspace{0.22in}
\namedateperiod
\vspace{0.18in}

\begin{learningtargetbox}
\small
\begin{enumerate}[leftmargin=1.5em, itemsep=5pt]
  \item \ldots{} read a \textbf{loss} $L(w)$ as a bowl and see that learning means finding its \textbf{lowest point} (the smallest error).
  \item \ldots{} use the \textbf{slope} to find downhill and take a gradient-descent step $w\leftarrow w-s\,L'(w)$, with $s$ the \textbf{learning rate}.
  \item \ldots{} extend the same rule to many weights with the \textbf{gradient} $\nabla L$, and explain how the learning rate controls the descent.
\end{enumerate}
\end{learningtargetbox}

\vspace{0.14in}

\begin{tocbox}
\small
\renewcommand{\arraystretch}{1.5}
\begin{tabularx}{\linewidth}{@{} c l X r @{}}
\rowcolor{blush}
\textbf{\#} & \textbf{Component} & \textbf{Description} & \textbf{Score} \\
\midrule
1 & Warm-Up        & Spiral review: squared-error loss, reading a slope for downhill, and one update step & \blank{1.2cm} \\
2 & Guided Notes   & The loss bowl; slope $=$ downhill; the update $w\leftarrow w-s\,L'(w)$; the gradient $\nabla L$ & \blank{1.2cm} \\
3 & Group Activity & Rolling downhill, descending from both sides, the learning rate, and the gradient of two weights & \blank{1.2cm} \\
4 & Exit Ticket    & Quick check: one step, when to stop (slope $=0$), and why we step against the slope & \blank{1.2cm} \\
5 & Homework       & Steps from both sides, the loss falling, learning-rate tuning, the gradient, and justifications & \blank{1.2cm} \\
\midrule
  & \multicolumn{2}{r}{\textbf{Total}} & \blank{1.2cm} \\
\end{tabularx}
\end{tocbox}

\vspace{0.10in}

\begin{remindbox}
\small We can \emph{build} networks --- layers, folds, filters --- but which \textbf{weights}? Every choice
gives a \textbf{loss}: the Unit-4 sum of squared errors between predictions and targets. Picture that loss as a
\textbf{bowl} over the weights; the best weights sit at the bottom. \textbf{Gradient descent} walks there: feel
the \textbf{slope} (it points uphill), step the opposite way $w\leftarrow w-s\,L'(w)$, and repeat until the
slope is $0$. The step size $s$ is the \textbf{learning rate}. With many weights the slope becomes the
\textbf{gradient} $\nabla L$, and the identical rule $\mathbf{w}\leftarrow\mathbf{w}-s\,\nabla L$ trains the
whole network.
\end{remindbox}

\end{document}
